% ============================================================
%  封面信息
% ============================================================
\school{计算机科学与技术学院}
\major{软件工程}
\student{2252441}{那晋宁}
\thesistitle{基于 RISC-V 和 Rust 的 Unix V6++操作系统内存管理、文件系统重构与移植}{}
\thesistitleeng{Refactoring and Porting Unix V6++ Memory Management and File System Based on RISC-V and Rust}{}
\thesisadvisor{方钰\quad 教授}
\thesisdate{2026}{5}{10}

% ============================================================
%  信息说明页
% ============================================================
\infotype{design}

\infoabstract{%
本设计面向同济大学 Unix V6++ 教学操作系统，重点完成内存管理、文件系统及相关 I/O 路径的 Rust 重构与 RISC-V 适配。针对页对象语义不清和释放责任分散的问题，设计重构物理页、页级分配、内核堆和交换区管理，引入类型化地址、Buddy 与 Slab 分层分配以及 RAII 资源回收。针对 inode、文件对象、文件描述符和缓冲区缓存之间引用关系分散的问题，设计在保留超级块、目录项和多级块映射的基础上，整理对象持有关系和错误传播路径。针对 RISC-V 迁移中的底层依赖变化，设计适配 Sv39 分页、MMIO UART 和 VirtIO 块设备。测试表明，Rust IA32 与 Rust RISC-V 版本均能完成文件系统装载、shell 运行、文件读写和用户程序执行。
}

% 毕业设计：图纸数量与正文字数
\infodrawings{0}
\infowordcount{约 21000}

\infomaterials{
  \item Unix V6++ Rust 重构与 RISC-V 移植源码；
  \item 测试程序与实验数据；
  \item 毕业设计论文及相关说明材料。
}
